\documentclass[14pt,usepdftitle=false,aspectratio=169]{beamer}
\usepackage{preambule}
\setbeamercolor{structure}{fg=black}
\usepackage{complexes,topologie}\def\addkern#1{\kern#1}
\hypersetup{pdftitle={Théorie spectrale -- \textit{C}\addkern{1pt}\textsuperscript{\textup{*}}-algèbres}}
\title{Théorie spectrale\\\emph{\textit{C}\addkern{1pt}\textsuperscript{\textup{*}}-algèbres}}
\author{}
\date{}
\begin{document}
\begin{frame}
    \titlepage
\end{frame}
\slideq{Propriété du spectre de $x$ autoadjoint dans une $C^*$-algèbre}{1/4}
\slider{$\sigma_A\l X\r\subset\bbR$}{1/4}
\slideq{$C^*$-algèbre}{2/4}
\slider{Algèbre de Banach munie d'une involution $\cdot^*\!:\!A\to A$ qui est antilinéaire ($\l\lambda x+\mu y\r^*=\bar\lambda x^*+\bar\mu y^*$), antimultiplicative ($\l xy\r^*=y^*x^*$) et telle que $\nrm{x^*x}=\nrm x^2$}{2/4}
\slideq{$\omega\l x^*\r$ pour $\omega\in\sigma\l A\r$}{3/4}
\slider{$\bar{\omega\l x\r}$}{3/4}
\slideq{Lien entre $\sigma_B$ et $\sigma_A$ si $B$ est une sous-$C^*$-algèbre de $A$}{4/4}
\slider{$\sigma_A\l x\r=\sigma_B\l x\r$ pour tout $x\in B$ car $B^\times=B\cap A^\times$}{4/4}
\end{document}